\begin{trapbox}
\textbf{易错一｜把分别递增写成全定义域递增。} 对 $f(x)=x-1/x$，$-1/2<1/2$，但 $f(-1/2)=3/2>f(1/2)=-3/2$。因此全定义域不是递增的。

\textbf{易错二｜混淆对勾函数。} $Ax+B/x$ 与 $Ax-B/x$ 的性质不同；后者每支值域均为实数，不能套 $2\sqrt{AB}$。

\textbf{易错三｜换元后忘记范围。} $t=x+1$ 时，必须先排除 $t=0$，再确定区间是否跨零。

\textbf{易错四｜把水平线两交点当成原点对称。} 两横坐标乘积固定，但只有 $c=0$ 时两根互为相反数。

\textbf{易错五｜忽略参数符号。} 本结论要求 $A,B>0$；若均为负，可整体提负号后得到每支递减；若异号则属于同号两项的对勾型，须另行分析。
\end{trapbox}
